\section{Artifact Availability and Limits}

The repository contains the \tool Python package, adapter scripts, metric
runners, generated CSV/Markdown/\LaTeX{} tables, provenance logs,
BibTeX/source-audit notes, an MIT license, and a reviewer quickstart. The entry
point is the tag
\texttt{audit-sid-cikm-\allowbreak resource-v0.1}, with a clean-checkout verifier for the
published tables and figure; the anonymous review artifact is available at
\url{https://anonymous.4open.science/r/sidinspector-9BB2}. Large local artifacts are kept out of git and
described through manifests so reviewers can separate public reproducibility
assets from local caches.

Users can apply \tool in three modes: downstream researchers flag incomplete,
aliased, or behaviorally weak artifacts before generator training; method
authors start from the minimal adapter template and receive D1--D5 reports; and
reviewers run the clean-check verifier against published CSVs. Runtime notes
separate lightweight D1/D2/D4/D5 preflight checks from the interaction-dependent
D3 step. These signals do not replace Recall@K or NDCG; they make tokenizer
artifacts inspectable before a full benchmark or generator-training run.

The evidence supports a resource-interface claim, not broad method coverage.
The GRID Musical row is a controlled feature-text export rather than an
end-to-end TIGER or GRID reproduction; ReSID is the bounded Musical export;
RQ-min is a local reference adapter; the Sports balanced GAOQ path was stopped
after constrained k-means became CPU-bound; and CARD remains a fallback/probe
path rather than faithful CARD reproduction~\cite{wei2026card}. D2 is not
causal collision harm; D3 uses prefix-neighborhood and reranker checks, not
trained-generator ranking; D5 is structural cost rather than measured latency;
D6 is optional churn evidence; and D7 requires generator traces absent from the
current rows.

New adapters emit the same normalized tables, declare their evidence role, and
record feature or checkpoint provenance. A named tokenizer improves coverage
only after item-level codes pass the validator and join to metadata or
interactions. A controlled mechanism probe improves diagnostic calibration but
does not count as named-method coverage; a generator trace activates D7.
Future versions should add causal collision-harm validation, broader faithful
adapter coverage, generator-output diagnostics, and unified search/recommendation
checks~\cite{li2024survey,penha2025jointsid}. Industrial SID deployment and
ranking studies further motivate treating these artifact properties as
engineering objects~\cite{singh2023bettersemanticids,li2026sidcoord,ju2026snapchatsid,xu2026largescalegenrec}.
